% ============ §14 Systems of parameters and multiplicity (原书 pp.104-110) ============
\section{Systems of parameters and multiplicity}\label{sec:14}\index{systems of parameters, section 14}\index{multiplicity, section 14}

Let $(A, m)$ be an $r$-dimensional Noetherian local ring; by Theorem~\ref{thm:13.4}, there exists an $m$-primary ideal generated by $r$ elements, but none
generated by fewer. If $a_1, \ldots, a_r \in m$ generate an $m$-primary ideal,
$\{a_1, \ldots, a_r\}$ is said to be a \emph{system of parameters}\index{system of parameters} of $A$ (sometimes
abbreviated to s.o.p.). If $M$ is a finite $A$-module with $\dim M = s$, there exist
$y_1, \ldots, y_s \in m$ such that $l(M/(y_1, \ldots, y_s)M) < \infty$, and then $\{y_1, \ldots, y_s\}$ is
said to be a \emph{system of parameters of $M$}.

If we set $A/m = k$, the smallest number of elements needed to generate
$m$ itself is equal to $\rank_k m/m^2$; (here $\rank_k$ is the rank of a free module
over $k$, that is the dimension of $m/m^2$ as $k$-vector space). This number is
called the \emph{embedding dimension}\index{embedding dimension} of $A$, and is written emb dim $A$. In general
\[
  \dim A \le \text{emb dim } A,
\]
and equality holds when $m$ can be generated by $r$ elements; in this case
$A$ is said to be a \emph{regular local ring}\index{regular local ring}, and a system of parameters generating
$m$ is called a \emph{regular system of parameters}.

\begin{theorem}\label{thm:14.1}
Let $(A, m)$ be a Noetherian local ring, and $x_1, \ldots, x_r$ a system
of parameters. Then
\begin{enumerate}[(i)]
  \item $\dim A/(x_1, \ldots, x_i) = r - i$ for $1 \le i \le r$.
  \item although it is not true that $\htop (x_1, \ldots, x_i) = i$ for all $i$ for an arbitrary
        system of parameters, there exists a choice of $x_1, \ldots, x_r$ such that every
        subset $F \subset \{x_1, \ldots, x_r\}$ generates an ideal of $A$ of height equal to the
        number of elements of $F$.
\end{enumerate}
\end{theorem}

\begin{proof}
(i) is contained in Theorem~\ref{thm:13.6}. We now prove the second half of
(ii). If $r \le 1$ the assertion is obvious; suppose that $r > 1$. Let $p_{0j}$ (for $1 \le
j \le e_0$) be the prime ideals of $A$ of height $0$. Choosing $x_1 \in m$ not contained
in any $p_{0j}$, we have $\htop (x_1) = 1$. Next letting $p_{1j}$ (for $1 \le j \le e_1$) be the minimal
prime divisors of $(x_1)$, so that $\htop p_{1j} = 1$, and choosing $x_2 \in m$ not contained
in any $p_{0j}$ or any $p_{1j}$, we have $\htop(x_1) = 1$, $\htop(x_1, x_2) = 2$; if $r = 2$ we're
done. If $r > 2$ we choose $x_3 \in m$ not contained in any minimal prime
divisor of $(0)$, $(x_1)$, $(x_2)$, $(x_1, x_2)$, and proceed in the same way to obtain
the result.

We now give an example where $\htop (x_1, \ldots, x_i) < i$. Let $k$ be a field and
set $R = k[X, Y, Z]$; let $I = (X) \cap (Y, Z)$, and write $A = R/I$, and $x, y, z$ for
the images in $A$ of $X, Y, Z$. The minimal prime ideals of $A$ are $(x)$ and
$(y, z)$; now $A/(x) = R/(X) \cong k[Y, Z]$ is two-dimensional and $A/(y, z) \cong
R/(Y, Z) \cong k[X]$ is one-dimensional, so that $\dim A = 2$. $\{y, x + z\}$ is a
system of parameters of $A$; in fact $xy = xz = 0$, so that $x^2 = x(x + z) \in
(y, x + z)$ and $z^2 = z(x + z) \in (y, x + z)$. However, $y$ is contained in the
minimal prime ideal $(y, z)$ of $A$, and hence $\htop (y) = 0$.
\end{proof}

\begin{theorem}\label{thm:14.2}
Let $(R, m)$ be an $n$-dimensional regular local ring, and
$x_1, \ldots, x_i$ elements of $m$. Then the following conditions are equivalent:
\begin{enumerate}[(1)]
  \item $x_1, \ldots, x_i$ is a subset of a regular system of parameters of $R$;
  \item the images in $m/m^2$ of $x_1, \ldots, x_i$ are linearly independent over $R/m$;
  \item $R/(x_1, \ldots, x_i)$ is an $(n - i)$-dimensional regular local ring.
\end{enumerate}
\end{theorem}

\begin{proof}
$(1) \Rightarrow (2)$ If $x_1, \ldots, x_i, x_{i+1}, \ldots, x_n$ is a regular system of parameters then
their images generate $m/m^2$ over $k = R/m$, and since $\rank_k m/m^2 = n$ they
must be linearly independent over $k$.

$(1) \Rightarrow (3)$ We know that $\dim R/(x_1, \ldots, x_i) = n - i$, and the images of
$x_{i+1}, \ldots, x_n$ generate the maximal ideal of $R/(x_1, \ldots, x_i)$.

$(3) \Rightarrow (1)$ If the maximal ideal $m/(x_1, \ldots, x_i)$ of $R/(x_1, \ldots, x_i)$ is generated
by the images of $y_1, \ldots, y_{n-i} \in m$ then $m$ is generated by $x_1, \ldots, x_i$,
$y_1, \ldots, y_{n-i}$.
\end{proof}

\begin{remark}
The hypothesis that $R$ is regular is not needed for $(3) \Rightarrow (1)$.
\end{remark}

\begin{proof}[Proof of $(2) \Rightarrow (1)$]
Using $\rank_k m/m^2 = n$, if we choose $x_{i+1}, \ldots, x_n \in m$ such that
the images of $x_1, \ldots, x_n$ in $m/m^2$ form a basis then $x_1, \ldots, x_n$ generate $m$
and so forms a regular system of parameters.
\end{proof}

\begin{theorem}\label{thm:14.3}
A regular local ring is an integral domain.
\end{theorem}

\begin{proof}
Let $(R, m)$ be an $n$-dimensional regular local ring; we proceed by
induction on $n$. If $n = 0$ then $m$ is an ideal generated by $0$ elements, so
that $m = (0)$. This in turn means that $R$ is a field. Thus a zero-dimensional
regular local ring is just a field by another name.

When $n = 1$, the maximal ideal $m = xR$ is principal and $\htop m = 1$, so that
there exists a prime ideal $p \neq m$ with $m \supset p$. If $y \in p$ we can write $y = xa$ with
$a \in R$, and since $x \notin p$ we have $a \in p$; hence $p = xp$, and by NAK, $p = (0)$. This
proves that $R$ is an integral domain. (There is a slightly different proof in the
course of the proof of Theorem~\ref{thm:11.2}; as proved there, a one-dimensional
regular local ring is just a DVR by another name.)

When $n > 1$, let $p_1, \ldots, p_l$ be the minimal prime ideals of $R$; then since
$m \nsubseteq m^2$ and $m \nsubseteq p_i$ for all $i$, there exists an element $x \in m$ not contained
in any of $m^2, p_1, \ldots, p_l$ (see Ex.~\ref{ex:1.6-ex}). Then the image of $x$ in $m/m^2$ is non-zero,
so that by the previous theorem $R/xR$ is an $(n - 1)$-dimensional regular
local ring. By the induction hypothesis, $R/xR$ is an integral domain, in
other words, $xR$ is a prime ideal of $R$. If $p_1$ is one of the minimal prime ideals
contained in $xR$ then since $x \notin p_1$, the same argument as in the $n = 1$ case
shows that $p_1 = xp_1$, and hence $p_1 = (0)$.
\end{proof}

\begin{theorem}\label{thm:14.4}
Let $(A, m, k)$ be a $d$-dimensional regular local ring; then
\[
  \gr_m(A) \cong k[X_1, \ldots, X_d],
\]
and if $\chi(n)$ is the Samuel function of $A$ then
\[
  \chi(n) = \binom{n+d}{d} \quad \text{for all } n \ge 0.
\]
\end{theorem}

\begin{proof}
Since $m$ is generated by $d$ elements, $\gr_m(A)$ is of the form
$k[X_1, \ldots, X_d]/I$, where $I$ is a homogeneous ideal. Now if $I \neq (0)$ let $S \in I$ be a
non-zero homogeneous element of degree $r$; then for $n \ge r$ the homogeneous
piece of $k[X]/I$ of degree $n$ has length at most
\[
  \binom{n+r-1}{d-1} - \binom{n-1}{d-1},
\]
which is a polynomial of degree $d - 2$ in $n$. This
implies that the Samuel function of $A$ is of degree at most $d - 1$, and
contradicts $\dim A = d$. Hence $I = (0)$; the second assertion follows from
the first.
\end{proof}

Let $(A, m)$ be a Noetherian local ring. Elements $y_1, \ldots, y_l \in m$ are said to be
\emph{analytically independent}\index{analytically independent elements} if they have the following property; for every
homogeneous form $F(Y_1, \ldots, Y_l)$ with coefficients in $A$,
\[
  F(y_1, \ldots, y_l) = 0 \Rightarrow \text{the coefficients of } F \text{ are in } m.
\]
If $y_1, \ldots, y_l$ are analytically independent and $A$ contains a field $k$, then
$F(y) \neq 0$ for any non-zero homogeneous form $F(Y) \in k[Y_1, \ldots, Y_l]$.

\begin{theorem}\label{thm:14.5}
Let $(A, m)$ be a $d$-dimensional Noetherian local ring and
$x_1, \ldots, x_d$ a system of parameters of $A$; then $x_1, \ldots, x_d$ are analytically
independent.
\end{theorem}

\begin{proof}
Set $q = \sum x_iA$. Since $q$ is an ideal of definition of $A$, by
Theorem~\ref{thm:13.4}, $\chi_q(n) = l(A/q^n)$ is a polynomial of degree $d$ in $n$ for $n \gg 0$. Set
$\bar{A} = k$; we say that a homogeneous form $\bar{f}(X) \in k[X_1, \ldots, X_d]$ of degree $n$
is a \emph{null-form} of $q$ if $F(x_1, \ldots, x_d) \in q^{n+1}m$ for any homogeneous form
$F(X) \in A[X_1, \ldots, X_d]$ which reduces to $\bar{f}(X)$ modulo $m$. Write $\mathfrak{n}$ for the ideal
of $k[X_1, \ldots, X_d]$ generated by the null-forms of $q$. Then
\[
  k[X_1, \ldots, X_d]/\mathfrak{n} = \bigoplus_{n \ge 0} q^n/q^{n+1}m,
\]
and writing $\varphi$ for the Hilbert polynomial of $k[X]/\mathfrak{n}$, we have $\varphi(n)
= l(q^n/q^{n+1}m)$ for $n \gg 0$. The right-hand side is just the number of elements in a
minimal basis of $q^n$, so that $\varphi(n) \cdot l(A/q) \ge l(q^n/q^{n+1})$. Now
\[
  l(q^n/q^{n+1}) = \chi_q(n) - \chi_q(n - 1)
\]
is a polynomial in $n$ of degree $d - 1$, so that $\deg \varphi \ge d - 1$, but if $\mathfrak{n} \neq (0)$ this
is impossible. Thus $\mathfrak{n} = (0)$, and the statement in the theorem follows at
once.
\end{proof}

\subsection*{Multiplicity}

Let $(A, m)$ be a $d$-dimensional Noetherian local ring, $M$ a finite $A$-module,
and $q$ an ideal of definition of $A$ (that is, an $m$-primary ideal). As we saw
in \S\ref{sec:13}, the Samuel function $l(M/q^{n+1}M) = \chi_M^q(n)$ can be expressed for
$n \gg 0$ as a polynomial in $n$ with rational coefficients, and degree equal to
$\dim M$, and therefore at most $d$. In addition, this polynomial can only take
integer values for $n \gg 0$, so it is easy to see by induction on $d$ (using the fact
that $\chi(n + 1) - \chi(n)$ has the same property) that
\[
  \chi_M^q(n) = \frac{e}{d!}n^d + \text{(terms of lower order)},
\]
with $e \in \Z$. This integer will be written $e(q, M)$. By definition we have the
following property.

\smallskip
\noindent\textbf{Formula 14.1.}\quad $\displaystyle e(q, M) = \lim_{n \to \infty} \frac{d!}{n^d}l(M/q^nM)$, and in particular, if $d = 0$ then
$e(q, M) = l(M)$.

\smallskip
\noindent From this we see easily the following:
\noindent\textbf{Formula 14.2.}\quad $e(q, M) > 0$ if $\dim M = d$, and $e(q, M) = 0$ if $\dim M < d$;

\smallskip
\noindent\textbf{Formula 14.3.}\quad $e(q^r, M) = e(q, M)r^d$;

\smallskip
\noindent\textbf{Formula 14.4.}\quad If $q$ and $q'$ are both $m$-primary ideals and $q \supset q'$ then
$e(q, M) \le e(q', M)$.

\smallskip
\noindent We set $e(q, A) = e(q)$, and define this to be the \emph{multiplicity} of $q$. In addition,
we will refer to the multiplicity $e(m)$ of the maximal ideal as the \emph{multiplicity of the
local ring} $A$, and sometimes write $e(A)$ for it. For example, if $A$ is a
regular local ring then by Theorem 4, we can see that $e(A) = 1$.

\begin{theorem}\label{thm:14.6}
Let $0 \to M' \to M \to M'' \to 0$ be an exact sequence of
finite $A$-modules. Then
\[
  e(q, M) = e(q, M') + e(q, M'').
\]
\end{theorem}

\begin{proof}
We view $M'$ as a submodule of $M$. Then
\[
  l(M/q^nM) = l(M''/q^nM'') + l(M'/M' \cap q^nM),
\]
and obviously $q^nM' \subset M' \cap q^nM$. On the other hand by Artin--Rees, there
exists $c > 0$ such that
\[
  M' \cap q^nM \subset q^{n-c}M' \quad \text{for all } n > c.
\]
Hence
\[
  l(M'/q^{n-c}M') \le l(M'/M' \cap q^nM) \le l(M'/q^nM').
\]
From this and Formula 14.1 it follows easily that
\[
  e(q, M) - e(q, M'') = \lim_{n \to \infty} \frac{d!}{n^d}l(M'/M' \cap q^nM) = e(q, M'). \qquad\blacksquare
\]
\end{proof}

\begin{theorem}\label{thm:14.7}
Let $\{p_1, \ldots, p_t\}$ be all the minimal prime ideals of $A$ such
that $\dim A/p_i = d$; then
\[
  e(q, M) = \sum_{i=1}^{t} l(M_{p_i})e(\bar{q}, A/p_i),
\]
where $\bar{q}_i$ denotes the image of $q$ in $A/p_i$, and $l(M_{p_i})$ stands for the length of $M_{p_i}$ as
$A_{p_i}$-module.
\end{theorem}

\begin{proof}[Proof (taken from Nagata [N1])]
We write $\sigma = \sum l(M_{p_i})$ and proceed by
induction on $\sigma$. If $\sigma = 0$ then $\dim M < d$, so that the left-hand side is $0$, and
the right-hand side is obviously $0$; now suppose $\sigma > 0$. Now there is some
$p \in \{p_1, \ldots, p_t\}$ for which $M_p \neq 0$; then $p$ is a minimal element of
$\Supp(M)$. Hence $p \in \Ass(M)$, that is $M$ contains a submodule $N$ isomorphic
to $A/p$. Then
\[
  e(q, M) = e(q, N) + e(q, M/N).
\]
On the other hand, $N_{p_i} \cong A_p/pA_p$, and $N_{p_i} = 0$ for $p_i \neq p$, so that $l(N_{p_i}) = 1$,
and the value of $\sigma$ for $M/N$ has decreased by one, so that the theorem holds
for $M/N$. However, from the definition
\[
  e(q, N) = e(q, A/p) = e(\bar{q}, A/p), \quad \text{where } \bar{q} = (q + p)/p.
\]
Putting this together, we see that the theorem also holds for $M$.
\end{proof}

Theorem 7 allows us to reduce the study of $e(q, M)$ to the case that $A$ is an
integral domain and $M = A$. In particular, if $A$ is an integral domain then
$l(M_{p_i})$ is just the rank of $M$, so that we obtain the following theorem.

\begin{theorem}\label{thm:14.8}
Let $A$ be a Noetherian local integral domain, $q$ an ideal of
definition of $A$ and $M$ a finite $A$-module; then
\[
  e(q, M) = e(q) \cdot s, \quad \text{where } s = \rank M.
\]
\end{theorem}

\begin{theorem}\label{thm:14.9}
Let $(A, m)$ be a Noetherian local ring, $q$ an ideal of definition
of $A$, and $x_1, \ldots, x_d$ a system of parameters of $A$ contained in $q$. Suppose that
$x_i \in q^{\nu_i}$ for $1 \le i \le d$. Then for a finite $A$-module $M$ and $s = 1, \ldots, d$ we have
\[
  e(q/(x_1, \ldots, x_s), M/(x_1, \ldots, x_s)M) \ge \nu_1\nu_2 \cdots \nu_s e(q, M).
\]
In particular if $s = d$, we have
\[
  l(M/(x_1, \ldots, x_d)M) \ge \nu_1\nu_2 \cdots \nu_de(q, M).
\]
\end{theorem}

\begin{proof}
It is enough to prove the case $s = 1$. We set $A' = A/x_1A$, $q' = q/x_1A$,
$M' = M/x_1M$ and $\nu = \nu_1$. By Theorem 1, we have $\dim A' = d - 1$. On the
other hand,
\[
  l(M'/q'^nM') = l(M/x_1M + q^nM)
\]
\[
  = l(M/q^nM) - l(x_1M + q^nM/q^nM).
\]
In addition, in view of $(x_1M + q^nM)/q^nM \cong x_1M/x_1M \cap q^nM \cong M/
(q^nM:x_1)$ and $q^{n-\nu}M \subset q^nM:x_1$, we have
\[
  l(M'/q'^nM') \ge l(M/q^nM) - l(M/q^{n-\nu}M),
\]
and therefore
\[
  l(M'/q'^nM') \ge l(M/q^nM) - l(M/q^{n-\nu}M).
\]
When $n \gg 0$ the right-hand side is of the form
\[
  \frac{e(q, M)}{d!}[n^d - (n - \nu)^d] + \text{(polynomial of degree } d - 2 \text{ in } n)
\]
\[
  = \frac{e(q, M)\nu}{(d-1)!}n^{d-1} + \text{(polynomial of degree } d - 2 \text{ in } n),
\]
so that the assertion is clear.
\end{proof}

A case of the above theorem which is particularly simple, but important,
is the following.

\begin{theorem}\label{thm:14.10}
Let $(A, m)$ be a $d$-dimensional Noetherian local ring, let
$x_1, \ldots, x_d$ be a system of parameters of $A$, and set $q = (x_1, \ldots, x_d)$; then
\[
  l(A/q) \ge e(q) \ge e(m).
\]
and if in addition $x_i \in m^{\nu_i}$ for all $i$ then $l(A/q) \ge \nu_1 \cdots \nu_d e(m)$.
\end{theorem}

\begin{theorem}\label{thm:14.11}
Let $A, m, x_i$ and $q$ be as above. Let $M$ be a finite $A$-module,
and set $A' = A/x_1A$, $M' = M/x_1M$ and $q' = q/x_1A = \sum_{i \ge 2} x_iA'$. Then if $x_1$ is a
non-zero-divisor of $M$, we have the following equality
\[
  e(q', M') = e(q, M).
\]
\end{theorem}

\begin{proof}
Since $l(M'/q'^{n+1}M') = l(M/x_1M + q^{n+1}M)$ we have
\[
  l(M/q^{n+1}M) - l(M'/q'^{n+1}M') = l(x_1M + q^{n+1}M/q^{n+1}M)
\]
\[
  = l(x_1M/x_1M \cap q^{n+1}M) = l(M/(q^{n+1}M:x_1))
\]
\[
  = l(M/q^nM) - l((q^{n+1}M:x_1)/q^nM).
\]
On the other hand, setting $\mathfrak{a} = \sum_{i \ge 2} x_iA$ we have $q = x_1A + \mathfrak{a}$ and $q^{n+1} =
x_1q^n + \mathfrak{a}^{n+1}$, and therefore
\[
  q^{n+1}M:x_1 = q^nM + (\mathfrak{a}^{n+1}M:x_1).
\]
Moreover, by Artin--Rees, there is a $c > 0$ such that for $n > c$ we have
$\mathfrak{a}^{n+1}M \cap x_1M = \mathfrak{a}^{n-c}(\mathfrak{a}^{c+1}M \cap x_1M)$, and therefore $\mathfrak{a}^{n+1}M:x_1 \subset \mathfrak{a}^{n-c}M$.
Thus
\[
  (q^{n+1}M:x_1)/q^nM = (q^nM + (\mathfrak{a}^{n+1}M:x_1))/q^nM
\]
\[
  \subset (q^nM + \mathfrak{a}^{n-c}M)/q^nM
\]
\[
  \cong \mathfrak{a}^{n-c}M/\mathfrak{a}^{n-c}M \cap q^nM.
\]
Now $\mathfrak{a}^{n-c}M/\mathfrak{a}^{n-c}M \cap q^nM$ is a module over $A/q'$, and since $\mathfrak{a}$ is generated
by $d - 1$ elements, $\mathfrak{a}^{n-c}$ is generated by $\binom{n-c+d-2}{d-2}$ elements. Thus
for $n > c$ we have
\[
  l(M'/q'^{n+1}M') \ge l(M/q^{n+1}M) - m\binom{n-c+d-2}{d-2},
\]
where $m$ is the number of generators of $M$. The right-hand side is a
polynomial of degree $d - 2$ in $n$, so that
\[
  e(q', M') = (d - 1)! \lim_{n \to \infty} l(M'/q'^{n+1}M')/n^{d-1}
\]
\[
  = (d - 1)! \lim_{n \to \infty} [l(M/q^{n+1}M) - l(M/q^nM)]/n^{d-1}
\]
\[
  = e(q, M).
\]
\end{proof}

\begin{theorem}[Lech's lemma]\label{thm:14.12}
Let $A$ be a $d$-dimensional Noetherian
local ring, and $x_1, \ldots, x_d$ a system of parameters of $A$; set $q = (x_1, \ldots, x_d)$, and
suppose that $M$ is a finite $A$-module. Then
\[
  e(q, M) = \lim_{\min(\nu_i) \to \infty} \frac{l(M/(x_1^{\nu_1}, \ldots, x_d^{\nu_d})M)}{\nu_1 \cdots \nu_d}.
\]
\end{theorem}

\begin{proof}
If $d = 0$ then both sides are equal to $l(M)$. If $d = 1$ then the right-hand
side is exactly Formula 14.1 which defines $e(q, M)$. For $d > 1$ we use
induction on $d$.

Setting $N_\nu = \{m \in M \mid x_1^\nu m = 0\}$ we have $N_1 \subset N_2 \subset \cdots$ so that there is a
$c > 0$ such that $N_\nu = N_{\nu+1} = \cdots$. If we set $M' = x_1^cM$ then $x_1$ is a non-zero-
divisor for $M'$, and there is an exact sequence $0 \to N_\nu \to M \to M' \to 0$.
Since $N_\nu$ is a module over $A/x_1^cA$ we have $\dim N_\nu < d$, and therefore
$e(q, N_\nu) = e(q, M')$. On the other hand,
\[
  l(M/(x_1^{\nu_1}, \ldots, x_d^{\nu_d})M) - l(M'/(x_1^{\nu_1}, \ldots, x_d^{\nu_d})M')
\]
\[
  = l(N_\nu + (x_1^{\nu_1}, \ldots, x_d^{\nu_d})M/(x_1^{\nu_1}, \ldots, x_d^{\nu_d})M)
\]
\[
  = l(N_\nu/N_\nu \cap (x_1^{\nu_1}, \ldots, x_d^{\nu_d})M)
\]
\[
  \le l(N_\nu/(x_1^{\nu_1}, \ldots, x_d^{\nu_d})N_\nu).
\]
If $\nu_1 \ge c$ then $x_1^{\nu_1}N_\nu = 0$, and $N_\nu$ is a module over the $(d - 1)$-dimensional
local ring $A/x_1^cA$, so that by induction there is a constant $C$ such that as
$\min(\nu_i) \to \infty$ we have
\[
  l(N_\nu/(x_1^{\nu_1}, \ldots, x_d^{\nu_d})N_\nu) = l(N_\nu/(x_2^{\nu_2}, \ldots, x_d^{\nu_d})N_\nu) \le C\nu_2 \cdots \nu_d.
\]
Therefore,
\[
  \lim [l(M/(x_1^{\nu_1}, \ldots, x_d^{\nu_d})M) - l(M'/(x_1^{\nu_1}, \ldots, x_d^{\nu_d})M')]/\nu_1 \cdots \nu_d = 0.
\]
This means that we can replace $M$ by $M'$ in the theorem, and so we can
assume that $x_1$ is a non-zero-divisor in $M$. Then by the previous theorem we
have $e(q, M) = e(\bar{q}, \bar{M})$, with $\bar{q} = q/x_1A$ and $\bar{M} = M/x_1M$. If we furthermore set
\[
  E = (x_2^{\nu_2}, \ldots, x_d^{\nu_d})M \quad \text{and} \quad F = M/E
\]
then by Theorem 9, we have
\[
  e(\bar{q}, \bar{M}) \cdot \nu_2 \cdots \nu_d \le l(M/(x_1^{\nu_1}, \ldots, x_d^{\nu_d})M) = l(F/x_1^{\nu_1}F)
\]
\[
  = \sum_{i=1}^{\nu_1} l(x_1^{i-1}F/x_1^iF) \le \nu_1l(F/x_1F) = \nu_1l(M/(x_1, x_2^{\nu_2}, \ldots, x_d^{\nu_d})M).
\]
Then by induction on $d$ we have
\[
  \lim l(M/(x_1^{\nu_1}, \ldots, x_d^{\nu_d})M)/\nu_1 \cdots \nu_d = \lim l(\bar{M}/(\bar{x}_2^{\nu_2}, \ldots, \bar{x}_d^{\nu_d})\bar{M})/\nu_2 \cdots \nu_d
  = e(\bar{q}, \bar{M}). \qquad\blacksquare
\]
\end{proof}

Although we will not use it in this book, we state here without proof a
remarkable result of Serre which shows that multiplicity can be expressed
as the Euler characteristic of the homology groups of the Koszul complex
(discussed in \S\ref{sec:16}).

\begin{theorem*}
Let $A$ be a $d$-dimensional Noetherian local ring, and $x_1, \ldots, x_d$ a
system of parameters of $A$; set $q = (x_1, \ldots, x_d)$ and let $M$ be a finite $A$-
module. Then
\[
  e(q, M) = \sum_i (-1)^i l(H_i(x, M)).
\]
For a proof, see for example Auslander and Buchsbaum [2].
\end{theorem*}

As we have seen in several of the above theorems, the multiplicity of
ideals generated by systems of parameters enjoy various nice properties. We
are now going to see that in a certain sense the general case can be reduced
to this one. We follow the method of Northcott and Rees [1].

More generally, let $A$ be a ring and $\mathfrak{a}$ an ideal. We say that an ideal $\mathfrak{b}$ is a
\emph{reduction}\index{reduction of an ideal} of $\mathfrak{a}$ if it satisfies the following condition:
\[
  \mathfrak{b} \subset \mathfrak{a}, \quad \text{and for some } r \ge 0 \text{ we have } \mathfrak{a}^{r+1} = \mathfrak{b}\mathfrak{a}^r.
\]
If $\mathfrak{b}$ is a reduction of $\mathfrak{a}$ and $\mathfrak{a}^{r+1} = \mathfrak{b}\mathfrak{a}^r$ then for any $n \ge 0$ we have $\mathfrak{a}^{r+n} = \mathfrak{b}^n\mathfrak{a}^r$.

\begin{theorem}\label{thm:14.13}
Let $(A, m)$ be a Noetherian local ring, $q$ an $m$-primary ideal
and $\mathfrak{b}$ a reduction of $q$; then $\mathfrak{b}$ is also $m$-primary, and for any finite $A$-module
$M$ we have
\[
  e(q, M) = e(\mathfrak{b}, M).
\]
\end{theorem}

\begin{proof}
If $q^{r+1} = \mathfrak{b}q^r$ then $q^{r+1} \subset \mathfrak{b} \subset q$, hence $\mathfrak{b}$ is also $m$-primary.
Moreover,
\[
  l(M/\mathfrak{b}^{n+1}M) \ge l(M/q^{n+1}M) = l(M/\mathfrak{b}^nq^r) \ge l(M/\mathfrak{b}^{n+r}M),
\]
so that $e(q, M) = e(\mathfrak{b}, M)$ follows easily.
\end{proof}

\begin{theorem}\label{thm:14.14}
Let $(A, m)$ be a $d$-dimensional Noetherian local ring, and
suppose that $A/m$ is an infinite field; let $q = (u_1, \ldots, u_s)$ be an $m$-primary
ideal. Then if $y_i = \sum a_{ij}u_j$ for $1 \le i \le d$ are $d$ `sufficiently general' linear
combinations of $u_1, \ldots, u_s$, the ideal $\mathfrak{b} = (y_1, \ldots, y_d)$ is a reduction of $q$ and
$\{y_1, \ldots, y_d\}$ is a system of parameters of $A$.
\end{theorem}

\begin{proof}
If $d = 0$ then $q^r = (0)$ for some $r > 0$, hence $(0)$ is a reduction of $q$
so that the result holds. We suppose below that $d > 0$.

\textbf{Step 1.} Set $A/m = k$ and consider the polynomial ring $k[X_1, \ldots, X_s]$ (or
$k[X]$ for short). For a homogeneous form $\varphi(X) = \varphi(X_1, \ldots, X_s) \in A[X]$
of degree $n$, we write $\bar{\varphi}(X) \in k[X]$ for the polynomial obtained by
reducing the coefficients of $\varphi$ modulo $m$. As in the proof of Theorem 5 we say that
$\bar{\varphi}(X) \in k[X]$ is a null-form of $q$ if $\varphi(u_1, \ldots, u_s) \in q^{n+1}m$; this notion depends
not just on $q$, but also on $u_1, \ldots, u_s$. However, for fixed $\bar{\varphi}$ it does not
depend on the choice of $\varphi$. We write $Q$ for the ideal of $k[X]$ generated by
all the null-forms of $q$, and call $Q$ the \emph{ideal of null-forms} of $q$. One sees
easily that all the homogeneous elements of $Q$ are null-forms of $q$, and
that the graded ring $k[X]/Q$ has graded component of degree $n$ isomorphic
to $q^n/q^{n+1}m$, so that we have
\[
  k[X]/Q = \bigoplus_{n \ge 0} q^n/q^{n+1}m = \gr_q(A) \otimes_A k.
\]
Write $\varphi(n)$ for the Hilbert function of $k[X]/Q$; then
\[
  l(q^n/q^{n+1}) \le \varphi(n) = l(q^n/q^{n+1}m) \le l(A/q) \cdot l(q^n/q^{n+1})
\]
(see the proof of Theorem 5). We know that for $n \gg 0$, the function $l(q^n/q^{n+1})$
is a polynomial in $n$ of degree $d - 1$ (where $d = \dim A$). Thus from the
above inequality, $\varphi$ is also a polynomial of degree $d - 1$, so that by
Theorem~\ref{thm:13.8}, (ii), we have $\dim k[X]/Q = d$.

Now set $V = \bigoplus kX_i$, and let $P_1, \ldots, P_t$ be the minimal prime divisors of
$Q$. By the assumption that $d > 0$, we have $P_i \nsupseteq V$, so that $P_i \cap V$ is a
proper vector subspace of $V$. Since $k$ is an infinite field,
\[
  V \neq \bigcup_{i=1}^{t} (V \cap P_i).
\]
Hence we can take a linear form $l_1(X) \in V$ not belonging to any $P_i$. If
$d \ge 1$ then similarly we can take $l_2(X) \in V$ such that $l_2(X)$ is not contained
in any minimal prime divisor of $(Q, l_1(X))$, and, proceeding in the same
way, we get $l_1(X), \ldots, l_d(X) \in V$ such that $(Q, l_1, \ldots, l_d)$ is a primary ideal
belonging to $(X_1, \ldots, X_s)$.

\textbf{Step 2.} We let $\mathfrak{b}$ be the ideal of $A$ generated by linear combinations
\[
  L_i(u) = \sum_{j=1}^{s} a_{ij}u_j \quad \text{(for } 1 \le i \le d\text{)}
\]
of $u_1, \ldots, u_s$ with coefficients in $A$. Then if we
put $l_i(X) = \sum a_{ij}X_j$, a necessary and sufficient condition for $\mathfrak{b}$ to be
a reduction of $q$ is that the ideal $(Q, l_1, \ldots, l_d)$ of $k[X]$ is $(X_1, \ldots, X_s)$-
primary.

\emph{Proof of necessity.} Suppose that $\mathfrak{b}q^r = q^{r+1}$. Then if $M(X)$ is a
monomial of degree $r + 1$ in $X_1, \ldots, X_s$, we can write
\[
  M(u) - \sum_i L_i(u)F_i(u) \in q^{r+2}m
\]
where the $F_i(X)$ are homogeneous forms of degrees $r$ with coefficients in
$A$. Thus
\[
  \bar{M}(X) - \sum_i l_i(X)\bar{F}_i(X) \in Q.
\]
Hence
\[
  (X_1, \ldots, X_s)^{r+1} \subset (Q, l_1, \ldots, l_d).
\]

\emph{Proof of sufficiency.} We go through the same argument in reverse: if
$\bar{M}(X) - \sum l_i\bar{F}_i \in Q$ then
\[
  M(u) - \sum L_i(u)F_i(u) \in q^{r+2}m,
\]
so that $q^{r+1} \subset \mathfrak{b}q^r + q^{r+2}m$; thus by NAK, $q^{r+1} = \mathfrak{b}q^r$.

\textbf{Step 3.} Putting together Steps 1 and 2 we see that $q$ has a reduction
$\mathfrak{b} = (y_1, \ldots, y_d)$ generated by $d$ elements. Both $q$ and its reduction $\mathfrak{b}$ are
$m$-primary ideals, so that $y_1, \ldots, y_d$ is a system of parameters of $A$. We
are going to prove that there exists a finite number of polynomials $D_\alpha(Z_{ij})$
for $1 \le \alpha \le \mu$ in $sd$ indeterminates $Z_{ij}$ (for $1 \le i \le d$ and $1 \le j \le s$) such that
$d$ linear combinations $y_i = \sum a_{ij}u_j$ (for $1 \le i \le d$) generate a reduction ideal
of $q$ if and only if at least one of $D_\alpha(a_{ij}) \neq 0$. (The expression `$d$ sufficiently
general linear combinations' in the statement of the theorem is quite vague,
but in the present case it has a precise interpretation as above.)

Let $G_1(X), \ldots, G_t(X)$ be a set of generators of $Q$, with $G_j$ homogeneous
of degree $e_j$. For any $sd$ elements $a_{ij}$ of $k$ (for $1 \le i \le d$ and $1 \le j \le s$), set
$l_i(X) = \sum a_{ij}X_j$. We write $I_n$ for the homogeneous component of degree $n$
of a homogeneous ideal $I \subset k[X_1, \ldots, X_s]$, and in particular we write
$(X_1, \ldots, X_s)_n = V_n$, so that
\[
  (X_1, \ldots, X_s)_n \subset (Q, l_1, \ldots, l_d) \iff V_n = (Q, l_1, \ldots, l_d)_n.
\]
Set $c_n = \dim_k V_n$. We have
\[
  (Q, l_1, \ldots, l_d)_n = \left\{\sum l_iF_i + \sum G_jH_j \mid F_i \in V_{n-1} \text{ and } H_j \in V_{n-e_j}\right\}.
\]
Let $K_1, \ldots, K_w$ be the elements obtained as $\sum l_iF_i + \sum G_jH_j$ as the $F_i$ run
through a basis of $V_{n-1}$ and the $H_j$ run independently through a basis
of $V_{n-e_j}$; it is clear that they span $(Q, l_1, \ldots, l_d)_n$. Each of $K_1, \ldots, K_w$ is a
linear combination of the $c_n$ monomials of degree $n$ in the $X_i$, with linear
functions in the $a_{ij}$ as coefficients; we write out these coefficients in a
$c_n \times w$ matrix. If $\psi_{\nu n}(a_{ij})$ for $1 \le \nu \le \rho$ are the $c_n \times c_n$ minors of this matrix
then the necessary and sufficient condition for $(X_1, \ldots, X_s)_n \subset (Q, l_1, \ldots, l_d)$ to
hold is that at least one of the $\psi_{\nu n}(a_{ij})$ is non-zero. Therefore the ideal
$(Q, l_1, \ldots, l_d)$ will fail to be $(X_1, \ldots, X_s)$-primary if and only if the
quantities $a_{ij}$ satisfy $\psi_{\nu n}(a_{ij}) = 0$ for all $n$ and all $\nu$. However, the ring $k[Z_{ij}]$
is Noetherian, so that the ideal of $k[Z_{ij}]$ generated by all of the $\psi_{\nu n}(Z_{ij})$ is
generated by finitely many elements $D_\alpha(Z_{ij})$ for $1 \le \alpha \le \mu$. These $D_\alpha$ clearly
meet our requirements.
\end{proof}

\begin{remark}
The polynomials $D_\alpha(Z_{ij})$ obtained above are in fact the necessary
and sufficient conditions on the coefficients $a_{ij}$ for the system of homo-
geneous equations $l_1(X) = \cdots = l_d(X) = G_1(X) = \cdots = G_t(X) = 0$ to have a
non-trivial solution, and as such they are known as a \emph{system of resultants}.
Here we have avoided appealing to the classical theory of resultants by
following a method given in Shafarevich [Sh].
\end{remark}

If $k = A/m$ is a finite field then Theorem 14 cannot to be used as it
stands, but we can use the following trick. Let $x$ be an indeterminate over
$A$, and set $S = A[x] - m[x]$; then $S$ consists of polynomials having a unit
of $A$ among their coefficients, and so the composite of the canonical maps
$A \to A[x] \to A[x]_S$ is injective. (In fact $S$ does not contain any zero-
divisors of $A[x]$, so that $A \subset A[x] \subset A[x]_S$; for this see [AM], Chap. 1,
Ex. 2.) Following Nagata~\cite{N1} we write $A(x)$ for $A[x]_S$. This is a
Noetherian local ring containing $A$, with maximal ideal $mA(x)$, and the
residue class field $A(x)/mA(x)$ is the field of fractions of $A[x]/m[x] = k[x]$,
that is, the field $k(x)$ of rational functions over $k$; this is an infinite field.
If $q$ is an $m$-primary ideal of $A$ then $qA(x)$ is a primary ideal belonging
to $mA(x)$. Moreover, since $A(x)$ is flat over $A$, we see that quite generally
if $I \supset I'$ are ideals of $A$ such that $I/I' \cong k$, then
\[
  IA(x)/I'A(x) \cong (I/I') \otimes_A A(x) \cong k \otimes_A A(x) = A(x)/mA(x).
\]
This gives $l_A(A/q^n) = l_{A(x)}(A(x)/q^nA(x))$, so that
\[
  \dim A = \dim A(x) \quad \text{and} \quad e(q) = e(qA(x)).
\]
Thus there are many instances when we can discuss properties of
$e(q)$ in terms of $A(x)$, to which Theorem 14 applies.

% ============ 习题 ============
\section*{Exercises to \S\ref{sec:14}}
\emph{Prove the following propositions.}

\begin{exercise}\label{ex:14.1-ex}
Let $(A, m)$ be a Noetherian local ring and set $G = \gr_m(A)$.
\begin{enumerate}[(i)]
  \item If $G$ is an integral domain then so is $A$ (hence Theorem 3 also follows
        from Theorem 4).
  \item Let $k$ be a field, and $A = k[[X, Y]]/(Y^2 - X^3)$; then $A$ is an integral
        domain, but $G$ has nilpotents.
\end{enumerate}
\end{exercise}

\begin{exercise}\label{ex:14.2-ex}
Let $(A, m)$ and $G$ be as above. For $a \in A$, suppose that $a \in m^i$ but $a \notin m^{i+1}$,
and write $a^*$ for the image of $a$ in $m^i/m^{i+1}$, viewed as an element of $G$;
define $a^*$ to be the \emph{leading term} of $a$. Set $0^* = 0$. Then
\begin{enumerate}[(i)]
  \item if $a^*b^* \neq 0$ then $(ab)^* = a^*b^*$;
  \item if $a^*$ and $b^*$ have the same degree and $a^* + b^* \neq 0$ then $(a + b)^* =
        a^* + b^*$;
  \item let $I \subset m$ be an ideal of $A$. Write $I^*$ for the ideal of $G$ generated by all
        the leading terms of elements of $I$; then setting $B = A/I$ and $n = m/I$, we
        have $\gr_n(B) = G/I^*$.
\end{enumerate}
\end{exercise}

\begin{exercise}\label{ex:14.3-ex}
In the above notation, if $G$ is an integral domain and $I = aA$ then
$I^* = a^*G$. If $I = (a_1, \ldots, a_r)$ with $r > 1$ then it can happen that
$I^* \neq (a_1^*, \ldots, a_r^*)$. Construct an example.
\end{exercise}

\begin{exercise}\label{ex:14.4-ex}
Let $(A, m)$ be a regular local ring, and $K$ its field of fractions.
\begin{enumerate}[(i)]
  \item For $0 \neq a \in A$, set $v(a) = i$ if $a \in m^i$ but $a \notin m^{i+1}$; then $v$ extends to an
        additive valuation of $K$.
  \item Let $R$ be the valuation ring of $v$; then $R$ is a DVR of $K$ dominating $A$.
        Let $x_1, \ldots, x_d$ be a regular system of parameters of $A$, and set
        $B = A[x_2/x_1, \ldots, x_d/x_1]$ and $P = x_1B$; then $P$ is a prime ideal of $B$ and
        $R = B_P$.
\end{enumerate}
\end{exercise}

\begin{exercise}\label{ex:14.5-ex}
In the above notation, if $0 \neq f \in m$ then $v(f)$ is equal to the multiplicity of
$A/(f)$.
\end{exercise}

\begin{exercise}\label{ex:14.6-ex}
(Associativity formula for multiplicities.) Let $A$ be a $d$-dimensional Noether-
ian local ring, $x_1, \ldots, x_d$ a system of parameters of $A$, $q = (x_1, \ldots, x_d)$, and
for $s \le d$ let $\mathfrak{a} = (x_1, \ldots, x_s)$. Write $\Gamma$ for the set of all prime divisors of $\mathfrak{a}$
satisfying $\htop p = s$, $\coht p = d - s$. Let $M$ be a finite $A$-module. Use Lech's
lemma to prove the following formula:
\[
  e(q, M) = \sum_{p \in \Gamma} e(q + p/p, M_p),
\]
(in particular, it follows that $\Gamma \neq \emptyset$).
\end{exercise}

\begin{remark}
The name of the formula comes from its connection with the associativity
of intersection product in algebraic geometry. For details, see [S3], pp. 84--5.
\end{remark}

\begin{exercise}\label{ex:14.7-ex}
Let $(A, m)$ be an $n$-dimensional Noetherian local integral domain, with
$n > 1$. If $0 \neq f \in A$ then $A_f$ is a Jacobson ring (see p. 34).
\end{exercise}
